\section{Receding-horizon execution and grounded feedback}
\label{sec:C10:feedback}
Open-loop planning selects a sequence before observing intermediate outcomes. Receding-horizon execution repeatedly plans, executes a prefix, observes the real result, and replans. The latter can correct state estimates and avoid compounding prediction errors, at the cost of additional decision and observation overhead. The correct horizon depends on action reversibility, environment volatility, model accuracy, and latency. A plan for a deterministic local calculation can be executed in a larger block than a sequence of irreversible operations against a changing external service.

ReAct interleaves reasoning and acting with environment observations, providing a concrete language-model-agent design for this feedback loop~\cite{yao2022react}. Its important boundary is that an action requests an external operation and an observation records a returned result. A model-generated continuation that imitates the observation is not equivalent. An execution harness should prevent these roles from being conflated in storage and training. A reasoning step can predict what a tool will return, but the prediction must remain distinguishable from the actual returned evidence.

Feedback has a source and scope. A compiler error identifies a problem relative to a language and configuration. A test failure identifies a violated assertion on an executed case. A human correction can clarify intent. A model critique supplies another statistical judgment. These signals should not all be assigned the same authority. A planner can use a weak critique to generate a hypothesis and then seek stronger evidence. Replacing an observed test result with a persuasive self-explanation reverses that hierarchy and can cause the agent to reason away a real failure.

Feedback should update the right state. A transient service failure may justify retrying without revising a general skill. A consistent schema mismatch may require changing the tool adapter. An incorrect remembered assumption may require invalidating a memory record. A failed plan can expose an error in the goal interpretation, world model, action implementation, or acceptance checker. If every failure produces only a new natural-language reflection, the system may accumulate explanations while leaving the actual defect unchanged. Retain diagnostic evidence sufficient to choose among these updates.

